\section{The actual native diagonal reduces to a full free quadratic complex}
\label{sec:136-diagonal-reduction}

\paragraph{Archive and scope.}
This addresses the nonlinear remaining condition for cote136,
physical p.~117/footer116/scan117 and physical
p.~121/footer120/scan121. We retain the actual native input
terms of ARCH-136-NATIVE-MIXED-TOR. The reduction below
identifies a specific nonlinear complex and its behavior under
the specified low-degree repair. It supplies neither a tail
annihilator nor a nonzero graph residue by itself.

For a homological graded module \(P\), write
\(\Gamma^2_{\mathrm{gr}}(P)\) for the arity-two part of its
free strictly graded-commutative PD algebra: even elements have
divided squares and odd elements have zero strict square.
Extend a linear differential as the special PD derivation.
This notation is distinct from applying an unproved derived
functor to a nonflat module.

\paragraph{ARCH-136-DIAGONAL-REDUCTION.}
Let \(P_\bullet\to L\) be a nonnegative free resolution,
and truncate it at homological degree \(r\geq2\), retaining
the actual last syzygy:
\[
 Q_i=P_i\ (i<r),\qquad
 Q_r=\operatorname{coker}(P_{r+1}\to P_r)=\Omega^rL,
 \qquad Q_i=0\ (i>r).
                                                               \tag{136.DR1}
\]
Then the natural map induces
\[
 H_n(\Gamma^2_{\mathrm{gr}}P)\xrightarrow{\sim}
 H_n(\Gamma^2_{\mathrm{gr}}Q)
                              \qquad(n\leq2r-2).
                                                               \tag{136.DR2}
\]
No flatness of \(Q_r\) is required.

Apply this with \(r=4\), \(L=S/(2,I_M)\), and the native
complex \(K\) of (136.NM2)--(136.NM4), reversing its indices.
The actual single-block auxiliary-arity-four component has
\[
 \begin{aligned}
 H^3&\cong H_5(\Gamma^2_{\mathrm{gr}}P),\\
 H^2&\cong H_6(\Gamma^2_{\mathrm{gr}}P).
 \end{aligned}                                                   \tag{136.DR3}
\]
These identifications hold integrally. The indicated groups
survive the auxiliary-graded RH1 repair of the single-block
simplicial model. Over \(\mathbb Z_{(2)}\), the identification
of the whole auxiliary-arity-four complex with
\(\Gamma^2_{\mathrm{gr}}K\) is an isomorphism in every form.

\begin{proof}
Put \(J=\ker(P\to Q)\). It is acyclic, is zero below
degree \(r\), and has
\[
        J_r=\operatorname{im}(P_{r+1}\to P_r),
                       \qquad J_i=P_i\quad(i>r).
\]
The map of strict quadratic complexes is degreewise onto,
since the free PD and exterior module functors preserve
surjections. Let \(B\) denote its kernel.
In total degree strictly below \(2r\), at most one factor
of a quadratic monomial can have degree at least \(r\).
There are therefore no diagonal terms with two truncated
factors in that range. Multiplication gives a canonical
identification
\[
                       B_n=(J\otimes P)_n\quad(n<2r).
                                                               \tag{136.DR4}
\]
Indeed, for degrees below \(r\) both factors are unchanged.
For a factor in degree \(r\), the other factor has degree
less than \(r\) and is free; tensoring the displayed kernel
\(J_r\) is exact. A factor of degree greater than \(r\)
maps to zero, again with only one such factor present.
These are all summands of the quadratic complex in that
range. The ordering signs agree with the graded tensor
differential. The differential of \(J_r\) is zero, so it
does not leave the specified summands.

The complex \(J\otimes P\) is acyclic: \(P\) is bounded
below and free, and the usual first-quadrant tensor filtration
has exact columns. If \(n\leq2r-2\), both a cycle in degree
\(n\) and a term bounding it in degree \(n+1\) lie in the
range of (136.DR4). Hence \(H_n(B)=0\) for every
\(n\leq2r-2\). The long exact sequence of the onto map
proves (136.DR2), including its last asserted degree. Nothing
has been asserted at degree \(2r-1\), where a bounding term
could meet the first quadratic bottom-syzygy contribution.

For the native application let \(a,b,e,z,K\) have the meaning
of (136.NM2), and let \(C\) denote just the graded coherent
parameter module. The old arity-two module is
\(Se\oplus Sz\), with \(e\) in form three and \(z\)
in form four; \(K=(Se\oplus Sz)\oplus C\) as a graded
module. Multiplication gives a cochain map
\[
 \Gamma^2_{\mathrm{gr}}K
            \longrightarrow (\text{actual native algebra})^{\text{arity }4}.
                                                               \tag{136.DR5}
\]
On \(C\)-quadratic terms and on the cross terms it is the
identity. The only remaining terms are
\[
 \gamma_2(z)\longmapsto3\gamma_4(b),\qquad
 ez\longmapsto3a\gamma_3(b),                            \tag{136.DR6}
\]
in forms eight and seven. The PD iteration and product
identities give precisely these coefficients. Thus (136.DR5)
is an isomorphism in every form over \(\mathbb Z_{(2)}\),
and is already an isomorphism in forms at most six over
\(\mathbb Z\). It therefore identifies the stated \(H^2\)
and \(H^3\) integrally. Reversal by four is an even shift,
so it preserves all graded signs; quadratic cochain degree
\(j\) is homological degree \(8-j\). Equation (136.DR2)
now gives (136.DR3).

We check the smaller auxiliary arities needed for RH1.
Arity one is the free unit disk, and arity two is \(K\),
which has no \(H^0\) or \(H^1\). For arity three, multiply
\(K\) by the arity-one unit complex \([a\to b]\).
Over \(\mathbb Z_{(2)}\), the resulting map onto the actual
arity-three component has precisely the kernel
\[
       S(ea)\xrightarrow{-1}S(eb-2za),
                 \qquad d(ea)=-(eb-2za).                \tag{136.DR7}
\]
The map on the remaining old form-six term is multiplication
by three, a unit there. Thus this is a true contractible
kernel and the arity-three component is acyclic locally.
Integrally, the same comparison is an isomorphism in the
low degrees relevant here; alternatively its form-one module
is the bottom parameter module times \(a\), whose differential
has the distinct injective component given by that same
module times \(b\). It has no form-zero part. Hence arity
three has zero \(H^0\) and \(H^1\) integrally as well.

Realize the native family on
\(\Delta[2]/\operatorname{sk}_1\Delta[2]\), as in the
mixed theorem. Multiple relative basis simplexes give direct
sums of the above low-arity components and tensor terms with
actual free unit contractions; these also have zero low
cohomology. The auxiliary-graded RH1 step is consequently an
isomorphism in arities below four. In arity four, new
zeroform generators and their zeroform identifying relations
cannot multiply a positive form without increasing arity,
since the auxiliary-zero part is just \(S\). Positive-form
modules and their differentials to higher positive forms
are unchanged there. Only new boundaries of one-forms appear.
The same is true at every iteration and at its filtered
colimit. This proves the asserted preservation of \(H^2\)
and \(H^3\), at the specified factorwise scope.
\end{proof}

\paragraph{The exact free-resolution independence needed here.}
Over the integer-torsion-free ring \(S\), the functor
\(H_n\Gamma^2_{\mathrm{gr}}\) on nonnegative free homological
complexes sends chain-homotopic maps to equal maps whenever
\(n\) is not divisible by four. In particular its groups in
degrees five and six are independent of the chosen free
resolution, by the ordinary comparison homotopies.

Here is a direct proof at exactly that scope. For a complex
\(C\), form \(C\oplus D\), where \(D\) is the direct sum
of the free unit disks on the bases of every \(C_i\), with
upper generator in degree \(i+1\) and lower generator in
degree \(i\). If \(h\) is a homotopy between maps \(g,f\),
the two usual endpoint maps into this cylinder differ by
\(dH+Hd\). Explicitly, write \(u_i(x)\) for the upper
generator on \(x\in C_i\) and \(v_i(x)=du_i(x)\) for
its lower partner. The endpoints are
\(i_0(x)=(x,0)\) and
\(i_1(x)=i_0(x)+v_i(x)+u_{i-1}(d_Cx)\), with the final
term omitted for \(i=0\). Evaluation sends the upper disk generator to
\(h(x)\), its lower partner to \(dh(x)\), and the old
complex by \(g\). Thus their evaluations are \(g,f\).
There is an actual decomposition
\[
 \Gamma^2_{\mathrm{gr}}(C\oplus D)=
 \Gamma^2_{\mathrm{gr}}C\ \oplus\ (C\otimes D)
                     \ \oplus\ \Gamma^2_{\mathrm{gr}}D.
                                                               \tag{136.DR10}
\]
The middle term contracts. In the last term, products of two
different unit disks also contract. For one disk in degrees
\(j\to j-1\), its divided square is contractible when
\(j\) is even; when \(j\) is odd it is the block
\(S\xrightarrow{2}S\) in degrees \(2j-1\to2j-2\).
Its only homology is \(S/2S\) in degree \(2j-2\), a
multiple of four. The upper kernel is zero because \(S\)
has no two-torsion. Therefore the cylinder projection induces
an isomorphism in every degree not divisible by four.
Both endpoint maps are its inverses in those degrees, proving
the assertion. Arbitrary ranks cause no problem, since all
quadratic monomials have finite support.

This applies to the full free resolution extending the actual
native truncation and to the explicit binary-digit resolution
with its \(\theta\) generator. It makes their \(H_5\)
and \(H_6\) canonically isomorphic by comparison maps,
preserving coefficients and weights. It does not assert
invariance of \(\Gamma\) under arbitrary quasi-isomorphisms,
nor invariance in the excluded homological degrees. This is
one precise bridge for transferring the norm calculation;
the strict minimal-retract bridge in
ARCH-136-QUADRATIC-TAIL-ACTION is another.

\paragraph{A uniform finite exponent for the computation.}
For the minimal full resolution used in
ARCH-136-ACTUAL-TAIL-SYZYGY, tensor with the degree-one
generator \(\theta\) satisfying \(d\theta=2\).
Left multiplication by \(\theta\) is a linear homotopy
\(dh+hd=2\operatorname{id}\). Its special derivation
extension to the strict quadratic complex gives
\[
                  dH+Hd=4\operatorname{id}
                         \quad\text{on }\Gamma^2_{\mathrm{gr}}P.
                                                               \tag{136.DR8}
\]
This proves a four-annihilator for these quadratic homology
groups. It gives no annihilator by an ordinary power of a
positive-weight coefficient tail.
The same four-annihilator holds for an arbitrary full free
resolution of \(S/(2,I_M)\): multiplication by two lifts the
zero map on the resolved module, so it is chain homotopic to
zero, and extending that homotopy as a special arity derivation
gives the same identity. No preferred multiplication on that
arbitrary resolution is needed.

\paragraph{Finite computational evidence.}
The saved program \texttt{diagonal\_square\_probe.py} constructs
integer matrices in each finite weight and checks \(d^2=0\)
exactly. It finds native diagonal \(H^3\) already at coefficient
weight five. The first order-two class is killed by every
tested positive digit coefficient, but testing all classes
in weight fifteen finds classes surviving the fresh coefficient
\(T^{[16]}\) in weight thirty-one. These finite tests do not
claim an unbounded nonannihilation theorem or a nonzero graph
residue. The exact coefficient-structure question is addressed
separately by the norm-sequence calculation.

The action test uses kernels modulo sixteen and boundaries
modulo four. Equation (136.DR8) and the natural universal
coefficient exact sequences identify
\[
 \operatorname{im}\bigl(H_5(C/16)\to H_5(C/4)\bigr)
                          =H_5(C)\subset H_5(C/4),      \tag{136.DR9}
\]
because the induced map on the adjacent torsion terms is
multiplication by four. Thus this finite test covers all
integral classes in its source weight. The modular image
algorithm was independently checked against 9,600 complete
small \(\mathbb Z/4\) image tests; the kernel algorithm
against 80 complete small \(\mathbb Z/16\) kernel tests.
These computational checks are not recorded as proof-checker
verification of the mathematical theorem.

\paragraph{Verification status.}
The author has read the complete saved proof, including the
last truncation boundary, exact native comparisons, RH1 scope,
the free-cylinder resolution bridge, and the modular coefficient
interpretation. The parent and reviewer15 completed whole-proof
readings of the reduction before the free-cylinder paragraph
was added and communicated PASS on 2026-11-22. The added
paragraph and arbitrary-resolution four-annihilator have also
received both reviewers' targeted whole-proof PASS. No computation
is being counted as a proof-checker verification.
